\documentclass[UTF8,11pt]{ctexart}
\usepackage[margin=2.1cm]{geometry}
\usepackage{booktabs}
\usepackage{longtable}
\usepackage{array}
\usepackage{enumitem}
\usepackage{xcolor}
\usepackage{hyperref}
\usepackage{amsmath}
\usepackage{amssymb}
\usepackage{titlesec}
\setCJKmainfont{SimSun}
\setCJKsansfont{Microsoft YaHei}
\setCJKmonofont{SimSun}
\setmainfont{Times New Roman}
\hypersetup{colorlinks=true,linkcolor=blue!50!black,urlcolor=blue!50!black,citecolor=blue!50!black}
\newcommand{\evidence}[1]{\textcolor{gray}{\footnotesize[#1]}}
\newcommand{\zh}{\par\noindent\textbf{中文.}\ }
\newcommand{\en}{\par\noindent\textbf{English.}\ }
\title{\bfseries 科研智能体的前沿方法与证据治理\\[2pt]
\large Frontier Methods and Evidence Governance for Research Agents\\[6pt]
\normalsize 调研、差距、改进路线 / Survey, Gaps, and Improvement Roadmap}
\author{Research Agent 项目组 / Research Agent Project}
\date{2026-10-25}

\begin{document}
\maketitle
\begin{abstract}
\noindent
本报告对当前科研智能体 pipeline 做静态闭环审计，并调研 2025--2026 年的论文、GitHub 开源代码与 PR。核心判断是：领域瓶颈正在从“能否自动生成论文”转向“论文中的 claim 能否被独立验证”。我们给出当前系统的差距清单、面向高质量文章与实验的改进设计，以及已经落地的三项加固：结构化 claim 数值重算、静态闭环审计、runtime job 持久化。本报告同时保留反例与不确定性，不把本地 demo 结果当作外部有效性证据。

\noindent\textbf{Abstract.} This report audits the current research-agent pipeline and surveys 2025--2026 papers, open-source code, and GitHub pull requests. The central finding is that the bottleneck is shifting from ``can an agent generate a paper?'' to ``can reviewers independently verify the claims in that paper?'' We present a gap analysis, a design for higher-quality manuscripts and experiments, and three implemented hardenings: structured numerical claim verification, a static closure audit, and persistent runtime jobs. We separate local prototype evidence from external validity and do not present demo results as scientific proof.
\end{abstract}

\tableofcontents
\newpage

\section{审计方法与证据等级 / Method and Evidence Levels}

\zh
本报告使用四类证据：(1) 对本仓库代码、SPEC、测试和产物做静态审计；(2) 对本地 \texttt{oss/}、\texttt{cortex-agent/} 与 \texttt{claude-code-main/} 做代码级对照；(3) 对 GitHub 代码与 PR 做抽样核查；(4) 对 arXiv 2025--2026 论文做摘要级调研。外部论文标注为“预印本/已核验元数据”，本地结论标注文件与行号；未做同任务复现实验的内容不作为效果结论。
\en
We use four evidence classes: (1) static audit of this repository's code, specs, tests, and artifacts; (2) code-level comparison with local \texttt{oss/}, \texttt{cortex-agent/}, and \texttt{claude-code-main/}; (3) sampled checks of GitHub code and pull requests; (4) abstract-level review of 2025--2026 arXiv papers. External papers are marked as preprint or metadata-verified; local claims carry file/line references. We do not treat anything without a same-task reproduction as an effect claim.

\subsection{当前审计的可复现结果 / Reproducible Audit Results}
\begin{longtable}{>{\raggedright\arraybackslash}p{0.27\linewidth} >{\raggedright\arraybackslash}p{0.18\linewidth} >{\raggedright\arraybackslash}p{0.45\linewidth}}
\toprule
场景 / Scenario & 结果 / Result & 解释 / Interpretation\\
\midrule
\endhead
真实 CC-BY 语料纵切 & 4 来源，2287 span，13 项审计 11 pass / 2 warn / 0 fail & 结构一致，但无 run/claim/closure；\texttt{structurally\_consistent\_but\_not\_closed}\\
Iris S5 纵切 & 修复 G2 语义后 13/13 pass & 发现并修复“设计 ready 被误当成 G2 已通过”\\
topic-only fixture & 8 个论文草稿章节，0 fail & 可读工作稿，不是完整论文\\
全量测试 & 150 tests OK & 工程一致性，不等于科学有效性\\
\bottomrule
\end{longtable}

\zh
静态审计命令为 \texttt{python -m research\_agent closure audit --bundle <dir>}，检查 span 来源、候选/设计引用、G2/protocol、RunCapsule/EA、claim draft、paper 槽位、schematic 图与 G5。它只检查结构，不重跑实验、不验证引用内容、不替代人签。
\en
The static audit command is \texttt{python -m research\_agent closure audit --bundle <dir>}. It checks span provenance, candidate/design references, G2/protocol, RunCapsule/EA, claim drafts, paper slots, schematic figures, and G5. It checks structure only; it does not rerun experiments, verify citation content, or replace human signatures.

\section{当前 Pipeline 的闭环判定 / Closure Verdict}

\begin{longtable}{>{\raggedright\arraybackslash}p{0.24\linewidth} >{\raggedright\arraybackslash}p{0.30\linewidth} >{\raggedright\arraybackslash}p{0.38\linewidth}}
\toprule
链路 / Link & 当前状态 / Current state & 说明 / Note\\
\midrule
\endhead
S1 $\rightarrow$ S2/3 & 真消费 span & 候选引用能回到 span，brief 尚未被充分利用\\
S2/3 $\rightarrow$ S4 & 部分 & 无 run\_dataset 时成立；有 run\_dataset 时被硬编码假设覆盖\\
S4 $\rightarrow$ S5 & 断开 & run 不读取 design/protocol 哈希；协议在 run 后冻结\\
S5 $\rightarrow$ S6 & 闭合 & 数据图与 EA 绑定\\
S5 $\rightarrow$ Paper & 仅 ID 绑定 & 未从真实 EA/claim 文件重算（本次已补 claim assertion）\\
Paper $\rightarrow$ Closure & 部分 & 槽位检查存在；claim 语义与引用内容仍未完全核验\\
Runtime $\rightarrow$ UI & 单次任务闭合 & job 已持久化；前端尚不能恢复轮询卡片\\
\bottomrule
\end{longtable}

\zh
结论：当前系统已经形成“S1 到可读工作稿”的证据档案闭环，以及“S5 到 EA、claim draft、closure”的实验闭合链；但两者还没有自动串成一条完整证据链。对外最稳妥的定位是“可审计的 dry-lab 研究运行时”，不是“通用 AI 科学家”。
\en
Verdict: the system has an evidence-archive loop from S1 to a readable working draft, and an experimental loop from S5 to EA, claim draft, and closure. The two loops are not yet automatically connected. The defensible positioning is an auditable dry-lab research runtime, not a general AI scientist.

\section{前沿方法图谱 / Frontier Landscape}

\begin{longtable}{>{\raggedright\arraybackslash}p{0.18\linewidth} >{\raggedright\arraybackslash}p{0.24\linewidth} >{\raggedright\arraybackslash}p{0.24\linewidth} >{\raggedright\arraybackslash}p{0.24\linewidth}}
\toprule
系统 / System & 核心机制 / Core mechanism & 强项 / Strength & 关键弱点 / Key weakness\\
\midrule
\endhead
AI Scientist v1/v2 & 端到端 idea--experiment--paper；v2 用 agentic tree search & 真正跑代码、自动评审、workshop 级产出 & 文献综述浅、novelty 误判、实验稳健性与引用问题\\
Agent Laboratory & 三阶段 + 每阶段 human feedback & 成本低、人在环、可产出代码与报告 & 自动评估与人工评估差距大\\
OpenScholar / PaperQA2 & 大规模全文检索 + 段落级引用 & 文献综合与 citation verification 强 & 不负责完整实验与 claim closure\\
Co-Scientist & generate--debate--evolve + tournament & 假设生成/排序强，HITL & 不是完整实验/证据链系统\\
ScientistOne / CoE & citation/\allowbreak numerical/\allowbreak method/\allowbreak conclusion claim + audit & 零幻觉引用、分数验证、方法--代码对齐 & 预印本；范围集中在可计算任务\\
MLReplicate & 6 个系统、45 篇稿件、人工+自动双评 & 揭示“自动评审高分但人工发现系统性缺陷” & 评估成本高；不提供系统设计\\
Verification Gap Survey & 35 篇系统、7 维审计 & 指出 83\% 放代码，只有 38\% 放种子/轨迹、38\% 做 novelty verification & 综述，不提出具体系统\\
RefVerifier / Citation on Point & proposition-level citation support & 证明“真实存在但错配”的引用难以检测 & 仍需检索/全文/人审\\
\bottomrule
\end{longtable}

\zh
两个数字值得牢记：MLReplicate 发现，自动评审接收的稿件中有 59\% 含伪造或无支撑 claim；验证缺口综述发现，24 个可运行系统中只有 38\% 发布种子或执行轨迹，只有 38\% 报告 novelty verification。这说明“论文外观质量”与“科学可验证性”之间存在系统性裂缝。
\en
Two numbers matter: MLReplicate found that 59\% of manuscripts accepted by automated review contained fabricated or unsupported claims; the verification-gap survey found that, among 24 runnable systems, only 38\% released seeds or execution traces and only 38\% reported any novelty-verification method. The gap between paper-like appearance and scientific verifiability is systemic.

\section{GitHub 代码与 PR 证据 / GitHub Code and PR Evidence}

\subsection{AI Scientist: provider fallback and writeup hardening}
\zh
PR \#135 增加 OpenAlex 作为 Semantic Scholar 的替代检索源，动机是 429 限流会让程序挂起；PR \#242 增加 Claude 4.x 支持，并在 writeup 中注入“只允许英文”的语言规则、提高 reflection 上限、修正 provider 前缀分派。这些 PR 说明：真实系统的失败往往不在“想法”，而在 provider、语言、LaTeX 编辑和重试边界。
\en
PR \#135 adds OpenAlex as an alternative to Semantic Scholar after 429 rate limits caused hangs. PR \#242 adds Claude 4.x support, injects an English-only language rule into writeup, raises the reflection limit, and fixes provider-prefix dispatch. The lesson: real failures are often provider, language, LaTeX-edit, and retry boundaries, not idea generation.

\subsection{PaperQA2: citation processing is still hard-coded}
\zh
PaperQA2 issue \#820 指出，citation prompt 可配置，但 citation processing 仍硬编码为 MLA 括号引用，调用者无法替换。这说明“引用生成”与“引用处理/核验”是两个不同层；我们的系统也必须把引用格式、引用支持和导出样式分开。
\en
PaperQA2 issue \#820 notes that the citation prompt is configurable but citation processing remains hard-coded to MLA parenthetical style. Citation generation, citation support, and export style are separate layers; our system should separate them too.

\subsection{本地开源对照 / Local open-source comparison}
\begin{itemize}[leftmargin=1.4em]
\item \textbf{AutoResearch}: 真跑子进程、run receipt、blind critic；但无文献/引用核验。
\item \textbf{AutoSci}: 完整工作流与 wiki/edge graph；但 \texttt{--auto} 可跳过人工门，evidence edge 仍偏字符串。
\item \textbf{AutoResearchClaw}: 三层 citation verification、数字白名单、HITL、checkpoint；最接近我们的竞品，但 claim verifier 尚未完全接线。
\item \textbf{RSI-Anything}: 真跑 GPU、hidden tests、flock；但无文献证据链。
\item \textbf{SoL-Pi}: byte-level quote + sha + 行号；是很好的 citation primitive，但不是科研流水线。
\end{itemize}

\section{关键设计不足 / Key Design Gaps}

\subsection{P0：影响科学有效性的缺口 / P0 gaps}
\begin{enumerate}[leftmargin=1.4em]
\item \textbf{全文入口}: 需要 PDF/TXT/MD + 页码/段落锚点 + 中文解析；否则任意真实课题没有证据。
\item \textbf{Claim 语义}: 需要 citation/numerical/method/conclusion 四类 claim，以及 supported/partial/contradicted/insufficient 的 support verdict；仅有自由文本 statement 不够。
\item \textbf{预注册顺序}: 应先冻结 ProtocolCard，再执行 run；当前 run 不消费 protocol/design 哈希。
\item \textbf{独立判定}: run 与 Judge 分离；需要 oracle/nop、clean rerun、FailureRecord 与 verdict。
\item \textbf{内容级引用支持}: 引用存在不等于支持该句；需要 proposition-level verification，尤其是 wrong-pinpoint 错误。
\item \textbf{身份与签名}: 需要 signer 身份、公钥签名、append-only 审计；仅有 human\_reason 不足以问责。
\end{enumerate}

\subsection{P1：影响文章质量的缺口 / P1 gaps}
\begin{enumerate}[leftmargin=1.4em]
\item \textbf{Section-level revision}: 全局大纲评审 + 局部修订，而不是整篇重新生成；LongCat/MIRA 的证据支持这一点。
\item \textbf{失败归因}: 单次 free-form reflection 不够；需要多假设 failure attribution（SAGE）。
\item \textbf{证据准入}: 只有独立验证过的 evidence 才能改写 persistence/principles（Triage-PiEvo）。
\item \textbf{上下文压缩}: 压缩时必须保留 claim/evidence ID、哈希与丢弃的 tool result 数。
\item \textbf{工具权限}: 检索/解析/run/写作分 read/write/network，拒绝在注册层做。
\item \textbf{前端双真相}: SIM 自动状态机仍可能与真实 runtime 混淆，需要隔离或移除。
\end{enumerate}

\section{面向高质量文章与实验的改进设计 / Design for High-Quality Manuscripts and Experiments}

\subsection{证据准入流水线 / Evidence admission pipeline}
\zh
建议把流程改成：SourceHit $\rightarrow$ 全文 $\rightarrow$ span（页码+字符区间）$\rightarrow$ claim card $\rightarrow$ support verdict $\rightarrow$ manuscript claim $\rightarrow$ slot 绑定 $\rightarrow$ closure。硬门包括：无全文不生成 span；无 support 不生成 result claim；无 EA 不写数字；无 claim 不写结果；无 closure 不导出。
\en
Proposed flow: SourceHit $\rightarrow$ full text $\rightarrow$ span (page + character offsets) $\rightarrow$ claim card $\rightarrow$ support verdict $\rightarrow$ manuscript claim $\rightarrow$ slot binding $\rightarrow$ closure. Hard gates: no full text, no span; no support, no result claim; no EA, no number; no claim, no results text; no closure, no export.

\subsection{结构化 Claim 与数值重算 / Structured claims and numeric recomputation}
\zh
本次已落地数值 assertion：每条 claim 可带 metric/comparator/value/tolerance/observed/status，\texttt{run claim} 在签署前从 metrics 重算，失败直接拒绝；closure 再次重算；paper 结果段只渲染 status=pass 的已验证数值。下一步要把 citation/method/conclusion 三类也结构化。
\en
We implemented numerical assertions: each claim can carry metric/comparator/value/tolerance/observed/status. \texttt{run claim} recomputes from metrics before signing and rejects mismatches; closure recomputes again; the paper Results section renders only verified values with status=pass. Next, citation, method, and conclusion claims should be structured too.

\subsection{预注册与独立判定 / Pre-registration and independent judging}
\zh
协议应先冻结，并绑定 design hash、dataset hash、metric、seeds、stopping rule；runner 不得修改。run 与 Judge 分离，Judge 只看冻结 rubric 与产物；输出 support/refute/inconclusive，并记录 FailureRecord。这样“跑出来”与“判定为支持”才是两件事。
\en
Freeze the protocol first, binding design hash, dataset hash, metric, seeds, and stopping rule; the runner cannot modify it. Separate run from Judge; the Judge sees only the frozen rubric and artifacts; output support/refute/inconclusive and a FailureRecord. ``It ran'' and ``it supports the claim'' must be two different events.

\subsection{Section-level 写作与修订 / Section-level writing and revision}
\zh
一次生成两万字正文会放大 drift。建议：ResearchSpec/大纲 $\rightarrow$ 章节契约 $\rightarrow$ 每节独立检索/证据 $\rightarrow$ 全局评审 $\rightarrow$ 局部修订。每节保留 claim/slot 账本，修订只改受影响段落，避免整篇重写导致引用漂移。
\en
One-shot generation of a 20k-word manuscript amplifies drift. Proposed: ResearchSpec/outline $\rightarrow$ section contracts $\rightarrow$ per-section retrieval/evidence $\rightarrow$ global review $\rightarrow$ local revision. Keep a claim/slot ledger per section; revision changes only affected passages to avoid citation drift.

\subsection{引用支持与新颖性核验 / Citation support and novelty verification}
\zh
引用核验应分三层：存在性（DOI/元数据）、上下文支持（句子是否被 span 支持）、跨文献一致性（是否遗漏反证）。RefVerifier 与 Citation-on-Point 说明 wrong-case 容易查，wrong-pinpoint 很难；因此需要全文定位与 proposition-level 判定，而不是只查 DOI。Novelty 不能只靠模型自评，需要最近工作检索 + 人工短名单 + ABANDON 理由。
\en
Citation verification should have three layers: existence (DOI/metadata), contextual support (does the span support the sentence?), and cross-paper consistency (are contradicting findings omitted?). RefVerifier and Citation-on-Point show wrong-case detection is easier than wrong-pinpoint detection; full-text localization and proposition-level verdicts are required. Novelty cannot rely on self-judgement; use nearest-work retrieval, human shortlists, and explicit ABANDON reasons.

\subsection{Claude Code 式运行时 / Claude Code-style runtime}
\zh
参考 \texttt{cortex-agent} 与本地 Claude Code 契约：单一 tool decision chain、read/write/network 权限、上下文压缩保留证据账本、blocking waitpoint、transcript/checkpoint/rewind、leaf subagent。不要照搬任意 Bash 与 bypassPermissions。本次已实现 job 持久化与静态审计；下一步是权限注册表、压缩与 waitpoint。
\en
Borrow from \texttt{cortex-agent} and the local Claude Code contract: a single tool decision chain, read/write/network permissions, compaction that preserves the evidence ledger, blocking waitpoints, transcript/checkpoint/rewind, and leaf subagents. Do not copy arbitrary Bash or bypassPermissions. We implemented job persistence and static audit; next are permission registry, compaction, and waitpoints.

\section{实施路线与验收 / Roadmap and Acceptance}

\begin{longtable}{>{\raggedright\arraybackslash}p{0.12\linewidth} >{\raggedright\arraybackslash}p{0.40\linewidth} >{\raggedright\arraybackslash}p{0.40\linewidth}}
\toprule
优先级 / Priority & 交付 / Deliverable & 验收 / Acceptance\\
\midrule
\endhead
P0.1 & 全文导入 + 页码/段落锚点 + 中文解析 & 20 个主题可定位到页/段；无全文不生成 span\\
P0.2 & structured claim: citation/numerical/method/conclusion & 计分数字 100\% 绑 EA；伪造/错配 0 通过\\
P0.3 & freeze-before-run + independent judge & protocol hash 进入 capsule；verdict 与人工复核一致率 \(\ge 90\%\)\\
P0.4 & 4 个 RWJT 任务 + human evaluation & oracle 通过；nop 低分；核心 claim 覆盖 \(\ge 90\%\)\\
P1.1 & section-level revision + global review & 未支持事实句占比 \(\le 10\%\)；引用漂移下降\\
P1.2 & failure attribution / evidence admission & failure 分类覆盖 100\%；无效 evidence 不进入原则层\\
P1.3 & permissions + compaction + waitpoint & 无绕过写 claim/gate；压缩后证据 ID 可恢复\\
\bottomrule
\end{longtable}

\zh
评测不能只用自动评审。MLReplicate 的结果表明，自动评审可能接受含伪造 claim 的稿件；需要人评、独立留出集、oracle/nop 与成本记录。建议同时报告：citation hallucination rate、citation support F1、numeric reproduction rate、method--code alignment、unsupported factual sentence rate、human quality score、cost/latency。
\en
Evaluation cannot rely on automated review alone. MLReplicate shows automated review can accept manuscripts with fabricated claims; human evaluation, held-out sets, oracle/nop, and cost accounting are required. Report citation hallucination rate, citation support F1, numeric reproduction rate, method--code alignment, unsupported factual sentence rate, human quality score, and cost/latency.

\section{结论 / Conclusion}

\zh
当前 pipeline 已经能跑通一条受治理的 dry-lab 纵切，并在本次加固后具备结构化 claim 数值重算、静态闭环审计与 runtime 任务恢复。但它距离“高质量文章/实验”的完整要求仍有明显差距，最大缺口是全文证据入口、内容级引用支持、独立判定和正式评测。我们的差异化应落在“本地优先、可审计、人工门的证据运行时”，而不是“通用 AI Scientist”或“效果 SOTA”。

\en
The pipeline already runs a governed dry-lab slice, and now includes structured numerical claim verification, static closure audit, and persistent runtime jobs. It is still far from a complete high-quality manuscript/experiment system; the largest gaps are full-text evidence ingestion, content-level citation support, independent judging, and formal evaluation. Our differentiation should be a local-first, auditable, human-gated evidence runtime, not a general AI scientist or an unverified SOTA claim.

\section*{参考文献 / References}
\begin{enumerate}[leftmargin=1.4em,itemsep=2pt]
\item Sakana AI. \emph{The AI Scientist}. arXiv:2408.06292; \emph{AI Scientist-v2}. arXiv:2504.08066.
\item Samuel Schmidgall et al. \emph{Agent Laboratory: Using LLM Agents as Research Assistants}. arXiv:2501.04227.
\item Akari Asai et al. \emph{OpenScholar: Synthesizing Scientific Literature with Retrieval-augmented LMs}. arXiv:2411.14199; Nature 2025.
\item Google DeepMind. \emph{AI co-scientist}. arXiv:2502.18864.
\item ScientistOne. \emph{Towards Human-Level Autonomous Research via Chain-of-Evidence}. arXiv:2605.26340.
\item \emph{Autonomous Research Agents: A Survey of AI Scientists and the Verification Gap}. arXiv:2608.05179.
\item \emph{MLReplicate: Benchmarking Autonomous Research Systems for Machine Learning Reproducibility}. arXiv:2605.16616.
\item \emph{Real PMIDs, Wrong Papers: Harness Design for Biomedical Claim-Verification Agents}. arXiv:2610.29812.
\item \emph{RefVerifier: Semi-Automated Reference Claim Verification for Scientific Manuscripts}. arXiv:2609.07652.
\item \emph{Is this Citation on Point?} arXiv:2608.12571.
\item \emph{One Reflection Is Not Enough: Self-Correcting Autonomous Research via Multi-Hypothesis Failure Attribution}. arXiv:2606.31478.
\item \emph{Triage-PiEvo: Evidence Admission for Principle-Evolving Scientific Agents}. arXiv:2610.14683.
\item \emph{IdeaScientist: Orchestrating Agents for Grounded Scientific Ideation}. arXiv:2610.04074.
\item \emph{SPRA: Structured Paper Representation for Agents}. arXiv:2610.21424.
\item \emph{LongCat-DeepResearch Technical Report}. arXiv:2609.36071.
\item GitHub: SakanaAI/AI-Scientist PR \#135 (OpenAlex support), PR \#242 (Claude 4.x / writeup hardening); Future-House/paper-qa issue \#820 (citation processing hook).
\item 本仓库：\texttt{research\_agent/pipeline.py}, \texttt{research\_agent/closure.py}, \texttt{research\_agent/runtime/jobs.py}, \texttt{research\_agent/modules/run.py}, \texttt{research\_agent/modules/paper.py}, \texttt{modules/*/SPEC.md}.
\end{enumerate}

\noindent\footnotesize
证据说明：外部论文以 arXiv 摘要/元数据为准，未逐篇复现；GitHub PR 以公开描述与 diff 为准；本地结论以代码、SPEC 与测试为准。所有链接访问日期为 2026-10-25。\\
Evidence note: external papers are assessed at abstract/metadata level and were not reproduced; GitHub PRs are assessed from public descriptions and diffs; local claims are grounded in code, specs, and tests. Links accessed 2026-10-25.

\end{document}
